\documentclass{beamer}

\usepackage{tikz}
\usetheme{Madrid}
\usecolortheme{beaver}

\usepackage{fontspec}
\setmonofont{JuliaMono-Regular.ttf}[Scale = 0.7]

\usepackage{amsmath,amssymb}
\usepackage{tikz-cd}
\usepackage{verbatim}
\usepackage{xcolor}

\title[Containers in Higher Kinds]
{Containers in Higher Kinds}
\author[Tian, Altenkirch, Gylterud]
{
  \textbf{Zhili Tian}\inst{1} 
  \and Thorsten Altenkirch\inst{1}\\
  \and Håkon Robbestad Gylterud\inst{2}
}

\institute[]
{
  \inst{1}%
  Functional Programming Lab\\
  School of Computer Science\\
  University of Nottingham
  \and
  \inst{2}%
  Department of Informatics\\
  University of Bergen
}

\date[TYPES 2026] % (optional)
{TYPES, May 2026}

\begin{document}
\frame{\titlepage}

\begin{frame}[fragile]{Old Story}

\begin{block}{Theorem}
Inductive types can be modeled as initial F-algebras.
\end{block}

\begin{exampleblock}{}
  \begin{columns}[c,onlytextwidth]
    \begin{column}{0.5\textwidth}
    \begin{verbatim}  
    data ℕ : Set where
      zero : ℕ
      suc  : ℕ → ℕ

    foldℕ : {A : Set}
      → (⊤ ⊎ A → A) → ℕ → A
    foldℕ α zero = α (inj₁ tt)
    foldℕ α (suc n) = α (inj₂ (foldℕ α n))
    \end{verbatim}
    \end{column}

    \begin{column}{0.5\textwidth}
    \begin{tikzcd}[column sep=huge]
      \texttt{⊤ ⊎ ℕ} \arrow[r , "\texttt{⊤⊎ (foldℕ α)}"] \arrow[d , "\texttt{[zero,suc]}"] & \texttt{⊤ ⊎ A} \arrow[d, "\texttt{α}"] \\
      \texttt{ℕ} \arrow[r , "\texttt{foldℕ α}"] & \texttt{A}
    \end{tikzcd}
    \end{column}
  \end{columns}
\end{exampleblock}

\end{frame}

\begin{frame}[fragile]{Old Story}

\begin{block}{Definition}
\textbf{Containers} correspond to polynomial functors:
\[ X \mapsto \sum_{s \in S} X^{P(s)} \]
\end{block}

\begin{exampleblock}{}
  \begin{columns}[T,onlytextwidth]
    \begin{column}{0.35\textwidth}
    \begin{verbatim}
    record Cont : Set₁ where
      constructor _◃_
      field
        S : Set
        P : S → Set
    \end{verbatim}
    \end{column}

    \begin{column}{0.60\textwidth}
    \begin{verbatim}
    ⟦_⟧ : Cont → Set → Set
    ⟦ S ◃ P ⟧ X = Σ[ s ∈ S ] (P s → X)

    ⟦_⟧₁ : (SP : Cont) → (X → Y) 
      → ⟦ SP ⟧ X → ⟦ SP ⟧ Y
    ⟦ SP ⟧₁ g (s , f) = s , g ∘ f
    \end{verbatim}
    \end{column}
  \end{columns}
\end{exampleblock}
\end{frame}

\begin{frame}[fragile]{Old Story}

\begin{block}{Definition}
\textbf{W-types} correspond to all strictly positive inductive types.
They are the initial algebras (least fixpoints) of containers.
\end{block}

\begin{exampleblock}{}
  \begin{columns}
    \begin{column}{0.5\textwidth}
    \begin{verbatim}  
    data W (SP : Cont) : Set where
      sup : ⟦ SP ⟧ (W SP) → W SP

    foldW : {SP : Cont} {X : Set}
      → (⟦ SP ⟧ X → X) → W SP → X
    foldW α (sup (s , f)) = α (s , foldW α ∘ f)
    \end{verbatim}
    \end{column}

    \begin{column}{0.5\textwidth}
      \begin{tikzcd}[column sep=huge]
        \texttt{⟦ SP ⟧ (W SP)} \arrow[r , "\texttt{⟦ SP ⟧₁ (foldW α)}"] \arrow[d , "\texttt{sup}"] & \texttt{⟦ SP ⟧ A} \arrow[d, "\texttt{α}"] \\
        \texttt{W SP} \arrow[r , "\texttt{foldW α}"] & \texttt{A}
      \end{tikzcd}
    \end{column}
  \end{columns}
\end{exampleblock}
\end{frame}

\begin{frame}[fragile]{More Types}

What about types like \texttt{List}?

\begin{exampleblock}{}
\begin{verbatim}
data List (A : Set) : Set where
  [] : List A
  _∷_ : A → List A → List A
\end{verbatim}
\end{exampleblock}

If we fix the type \texttt{A}, then signature is:

\begin{exampleblock}{}
\begin{verbatim}
ListASig : (A : Set) → Set → Set
ListASig A X = ⊤ ⊎ A × X
\end{verbatim}
\end{exampleblock}
\end{frame}

\begin{frame}[fragile]{More Types}

What about types like \texttt{Bush}? (legit Agda definition btw)

\begin{exampleblock}{}
\begin{verbatim}
data Bush (A : Set) : Set where
  [] : Bush A
  _∷_ : A → Bush (Bush A) → Bush A
\end{verbatim}
\end{exampleblock}

If we fix the type \texttt{A}, then signature is:

\begin{exampleblock}{}
\begin{verbatim}
BushASig : (A : Set) → Set → Set
BushASig A X = ⊤ ⊎ A × ???
\end{verbatim}
\end{exampleblock}

\pause
Let's try to take \texttt{Bush} and \texttt{A} separately:

\begin{exampleblock}{}
\begin{verbatim}
BushSig : (Set → Set) → Set → Set
BushSig F X = ⊤ ⊎ F (F X) × X
\end{verbatim}
\end{exampleblock}
\end{frame}

\begin{frame}[fragile]{2nd-Order Functors}

\begin{block}{Definition}
A \textbf{2nd-order functor} has all following constructions:
\end{block}

\pause
\begin{exampleblock}{}
\begin{verbatim}
H : (Set → Set) → Set → Set
H F X = ⊤ ⊎ F (F X) × X
\end{verbatim}
\end{exampleblock}

\pause
\begin{exampleblock}{}
\begin{verbatim}
ℍ : Func → Func
ℍ (F , F₁) = H F , HF₁ where
  HF₁ : (X → Y) → H F X → H F Y
  HF₁ = ...

ℍ₁ : {𝔽 𝔾 : Func} → Nat 𝔽 𝔾 → Nat (ℍ 𝔽) (ℍ 𝔾)
ℍ₁ = ...
\end{verbatim}
\end{exampleblock}
\end{frame}

\begin{frame}[fragile]{2nd-Order Containers}

\begin{block}{Definition}
A \textbf{2nd-order container} is given by a shape, positions of both \texttt{X} and \texttt{F},
and recurively, other 2nd-order container for each position of \texttt{F}.
\end{block}

\begin{exampleblock}{}
\begin{verbatim}
record 2Cont : Set₁ where
  pattern
  inductive
  field
    S : Set
    PX : S → Set
    PF : S → Set
    RF : (s : S) → PF s → 2Cont
\end{verbatim}
\end{exampleblock}
\end{frame}

\begin{frame}[fragile]{2nd-Order Containers to Functors}

Function interpretation:

\begin{exampleblock}{}
\begin{verbatim}
2⟦_⟧ : 2Cont → (Set → Set) → Set → Set
2⟦ S ◃ PX + PF + RF ⟧ F X = 
  Σ[ s ∈ S ] ((PX s → X) × ((pF : PF s) → F (2⟦ RF s pF ⟧ F X)))
\end{verbatim}
\end{exampleblock}

\pause
Functor of $\textbf{\underline{[Set,Set]}}$ interpretation:

\begin{exampleblock}{}
\begin{verbatim}
2⟦_⟧F : 2Cont → Func → Func
2⟦_⟧F₁ : 2Cont → {𝔽 𝔾 : Func} → Nat 𝔽 𝔾 → Nat (ℍ 𝔽) (ℍ 𝔾)
\end{verbatim}
\end{exampleblock}
\end{frame}

\begin{frame}[fragile]{Initial 2nd-Order Algbra}

\begin{block}{Definition}
\textbf{2nd-order W-types} correspond to all strictly positive inductive types with kinds \texttt{(Set → Set)}. They are the initial algebras (least fixpoint) of 2nd-order containers.
\end{block}

\begin{exampleblock}{}
  \begin{columns}
    \begin{column}{0.5\textwidth}
    \begin{verbatim}  
    2W : 2Cont → Func

    2sup : {H : 2Cont} 
      → NatTrans (2⟦ H ⟧F (2W H)) (2W H)

    fold2W : {𝔽 : Func} {H : 2Cont} 
      → NatTrans (2⟦ H ⟧ 𝔽) 𝔽
      → NatTrans (2W H) 𝔽
    \end{verbatim}
    \end{column}

    \begin{column}{0.5\textwidth}
      \begin{tikzcd}[column sep=huge]
        \texttt{2⟦ H ⟧F (2W H)} \arrow[r , "\texttt{2⟦ H ⟧F₁ (fold2W α)}"] \arrow[d , "\texttt{2sup}"] & \texttt{2⟦ H ⟧F 𝔽} \arrow[d, "\texttt{α}"] \\
        \texttt{2W H} \arrow[r , "\texttt{fold2W α}"] & \texttt{𝔽}
      \end{tikzcd}
    \end{column}
  \end{columns}
\end{exampleblock}
\end{frame}

\begin{frame}[fragile]{More Types}

What about types like monad transformers?

\begin{exampleblock}{}
\begin{verbatim}
record MaybeT (M : Set → Set) (X : Set) : Set where
  field
    runMaybeT : M (Maybe X)
\end{verbatim}
\end{exampleblock}

Its signature is in even higher kinds:

\begin{exampleblock}{}
\begin{verbatim}
SigMaybeT : ((Set → Set) → Set → Set) → (Set → Set) → Set → Set
SigMaybeT _ M X = M (⊤ ⊎ X)
\end{verbatim}
\end{exampleblock}
\end{frame}

\begin{frame}[fragile]{Higher Kinds}

Syntax and semantics of kinds:

\begin{exampleblock}{}
\begin{verbatim}
data Kind : Set where
  * : Kind
  _⇒_ : Kind → Kind → Kind

⟦_⟧K : Kind → Set₁
⟦ * ⟧K = Set
⟦ A ⇒ B ⟧K = ⟦ A ⟧K → ⟦ B ⟧K
\end{verbatim}
\end{exampleblock}

\end{frame}

\begin{frame}[fragile]{Hereditary Functors}

\textbf{Hereditary functors} and \textbf{hereditary categories} are mutually defined:

\begin{exampleblock}{}
\begin{verbatim}
isHFunc : ⟦ A ⟧K → Set₁
isHFunc = ...

HFunc : Kind → Set₁
HFunc A = Σ (⟦ A ⟧K) isHFunc

HCat : Kind → Cat
HCat A = record { Obj = HFunc A ; Mor = NatTrans ; ... }
\end{verbatim}
\end{exampleblock}

\pause
\begin{alertblock}{Note}
Hereditary categories are not freely generated functor categories!
\end{alertblock}

\end{frame}

\begin{frame}[fragile]{Hereditary Containers}

Unnormalized syntax ($\lambda$-terms + polynomials)

\begin{exampleblock}{}
\begin{verbatim}
data Tm : Con → Kind → Set₁ where
  var : Var Γ A → Tm Γ A
  lam : Tm (Γ ▹ A) B → Tm Γ (A ⇒ B)
  app : Tm Γ (A ⇒ B) → Tm Γ A → Tm Γ B
  Π : (I : Set) → (I → Tm Γ A) → Tm Γ A
  Σ : (I : Set) → (I → Tm Γ A) → Tm Γ A
\end{verbatim}
\end{exampleblock}
\end{frame}

\begin{frame}[fragile]{Normalization}

Normalizing function:

\begin{exampleblock}{}
\begin{verbatim}
nf : Tm Γ A → Nf Γ A
nf = ...
\end{verbatim}
\end{exampleblock}

\begin{itemize}
  \item hereditary substitution:\\
    \begin{center}
    Recursive $\beta$-reduction (substitutions) on structures
    \end{center}
  \item container normalization:
    \[ \Sigma_{x:A} \Sigma_{y:B(x)} C(x,y) \cong \Sigma_{(x,y):\Sigma_{x:A} B(x)} C(x,y) \]
    \[ \Pi_{x:A} \Pi_{y:B(x)} C(x,y) \cong \Pi_{(x,y):\Sigma_{x:A} B(x)} C(x,y) \]
    \[ \Pi_{x:A} \Sigma_{y:B(x)} C(x,y) \cong \Sigma_{f:\Pi_{x:A} B(x)} \Pi_{x:A} C(x, f(x))\]
\end{itemize}
\end{frame}

\begin{frame}[fragile]{Hereditary Containers}

Normal forms, neutral terms and spines are mutually defined:

\begin{exampleblock}{}
\begin{verbatim}
data Nf : Con → Kind → Set₁ where
  ...

record Ne (Γ : Con) (B : Kind) : Set₁ where
  inductive
  field
    S : Set
    P : Var Γ A → S → Set
    R : (x : Var Γ A) (s : S) → P x s → Sp Γ A B

data Sp : Con → Kind → Kind → Set₁ where
  ...
\end{verbatim}
\end{exampleblock}
\end{frame}

\begin{frame}[fragile]{Hereditary Containers}

\begin{block}{Definition}
\textbf{Hereditary Containers} are closed normal forms
\end{block}

\begin{exampleblock}{}
\begin{verbatim}
HCont : Kind → Set₁
HCont A = Nf ∙ A
\end{verbatim}
\end{exampleblock}

\pause
Function interpretation:

\begin{exampleblock}{}
\begin{verbatim}
⟦_⟧ : HCont A → ⟦ A ⟧K
⟦_⟧ = ...
\end{verbatim}
\end{exampleblock}

\pause
Full semantics?

\begin{exampleblock}{}
\begin{verbatim}
⟦_⟧ : HCont A → HFunc A
⟦_⟧ = {- TODO -}
\end{verbatim}
\end{exampleblock}
\end{frame}

\begin{frame}{Conclusion}

\centering
\begin{tabular}{|c|c|c|}
  \hline
  \texttt{Bool, ℕ} & $\mu$ of Func & W on Cont \\
  \hline
  \texttt{ℕ∞} & $\nu$ of Func & M on Cont \\
  \hline
  \texttt{Fin} & $\mu$ of IFunc & IW on ICont \\
  \hline
  \texttt{Fin∞} & $\nu$ of IFunc & IM on ICont \\
  \hline
  \texttt{List, Bush} & $\mu$ of 2Func(\checkmark) & 2W(\checkmark) on 2Cont(\checkmark) \\
  \hline
  \texttt{Stream} & $\nu$ of 2Func & 2M(?) on 2Cont \\
  \hline
  \texttt{MaybeT, StateT} & $\mu$ of HFunc(\checkmark) & HW(?) on HCont(\checkmark) \\
  \hline
  \texttt{Vec : Set → ℕ → Set} & $\mu$ of (?) & (?) on (?) \\
  \hline
  ... & ... & ... \\
\end{tabular}

\begin{itemize}
  \item Fill in all (?) in the table.
  \item Prove all types reduce to normal W-types.
  \item Look at mutually inductive and coinductive types.
\end{itemize}

\end{frame}

\begin{frame}{}
\Huge
\[\int_c Talk\]
\end{frame}

\begin{frame}{}

Citation

\end{frame}

\end{document}